\chapter{Случайные процессы}

\begin{definition}
    \emph{Случайный процесс} $ \xi_t $ представляет собой заданное на вероятностном пространстве
    однопараметрическое семейство случайных величин с параметром $ t \in T $ (в качестве $ T $ можно
    брать различные множества).
\end{definition}

$ \xi_t(\omega) $ для $ \omega \in \Omega $ можно рассматривать
\begin{itemize}
    \item при фиксированном значении $ t $ как значения случайной величины;
    \item при фиксированных $ \omega $ как траекторию случайного процесса.
\end{itemize}

\begin{definition}
    Для любого $ n $ и при любых значениях $ t(1), t(2), \dotsc \in T $ функции
    $$ F_{t(1), \dots, t(n)}(x_1, \dots, x_n) = P\Set{\xi_{t(1)} < x_1, \quad \dots, \quad
    \xi_{t(n)} < x_n} $$
    представляют собой \emph{конечномерные} распределения.
\end{definition}

\begin{notation}
    Введём обозначения:
    \begin{itemize}
        \item \emph{Математическое ожидание}: $ \mathcal A(t) = \mathbb E \xi_t $.
        \item \emph{Дисперсия}: $ \mathcal D(t) = \mathcal D \xi_t $.
        \item \emph{Ковариационная (корреляционная) функция}: $ \mathcal B(t, s) =
            \mathcal B(\xi_t, \xi_s) = \mathbb E \xi_t \xi_s $.
    \end{itemize}
\end{notation}

\begin{definition}
    Процесс $ \xi_t $ называется \emph{процессом с независимыми приращениями}, если
    $$ \forall n = 2, 3, \dotsc \quad \forall t(1) < \dots < t(n) \quad
    \text{независимы приращения} \xi_{t(1)}, \quad \xi_{t(2)} - \xi_{t(1)}, \quad \dots, \quad
    \xi_{t(n)} - \xi_{t(n - 1)} $$
\end{definition}

\begin{remark}
    Для таких процессов достаточно знать их одномерные распределения.
\end{remark}

\begin{definition}
    Процесс с независимыми приращениями имеет \emph{однородные (стационарные)} приращения, если
    для любых $ t $ и $ \tau > 0 $ распределения разностей $ \xi_{t + \tau} - \xi_{t} $ не
    зависят от $ t $.
\end{definition}

\begin{remark}
    Для таких процессов верно
    $$ \mathcal A(t) = \mathcal A(s) + \mathcal A(t - s) $$
    $$ \mathcal D(t) = \mathcal D(s) + \mathcal D(t - s), $$
    а значит, $ \mathcal A $ и $ \mathcal D $ являются линейными функциями.
\end{remark}

\begin{definition}
    Процесс называется \emph{стационарным в узком смысле}, если для любого $ n $ и любых
    $ t(1) < \dots < t(n) $ многомерные распределения случайных величин
    $$ \xi_{t(1) + \tau)}, \quad \dots, \quad \xi_{t(1) + \tau} $$
    не зависят от $ \tau $.
\end{definition}

\begin{remark}
    Для таких процессов выполняется
    $$ \mathcal A(t) = \const, \quad \mathcal D(t) = \const, \quad
    \mathcal B(t, s) = \mathcal B(s - t) $$
\end{remark}

\begin{definition}
    Процесс называется \emph{стационарным в широком смысле}, если
    $$ \mathcal A(t) = \const, \quad \mathcal B(t, s) = \mathcal B(s - t) $$
\end{definition}

\begin{exmpls}
\item Любая последовательность независимых одинаково распределённых случайных величин
    является стационарным в узком смысле процессом.
\item Последовательность независимых случайных величин, у которых совпадают математические
    ожидания и дисперсии, представляет собой стационарный в широком смысле процесс.
\end{exmpls}

\section{Пуассоновские процессы}

\begin{definition}
    Процесс с независимыми однородными приращениями называется \emph{пуассоновским} с параметром
    $ \lambda > 0 $, если $ \xi_0 = 0 $, а $ \xi_t $ при каждом $ t > 0 $ имеет пуассоновское
    распределение с параметром $ \lambda t $.
\end{definition}

Траектории такого процесса имеют единичные скачки в случайных точках вида $ z_1 + \dots z_k $, где
$ z_i $ независимы и имеют экспоненциальное распределение с параметром $ \frac1\lambda $.

\section{Винеровские процессы}

\begin{definition}
    Процесс $ W(t) $ с независимыми однородными приращениями называется \emph{винеровским},
    если $ W(0) = 0 $, а случайная величина $ W(t) $ при любом фиксированном $ t > 0 $ имеет
    нормальное распределение $ N(0, t) $.
\end{definition}

\begin{undefthm}{Моментные характеристики винеровского процесса}
    $$ \mathbb E W(t) = 0, \quad \mathcal D W(t) = t $$

    Пусть $ s \le t $.
    \begin{multline*}
        \mathcal B(s, t) = \mathbb EW(s)W(t) = \mathbb E W(s) \Bigl( W(s) + \bigl( W(t) -
        W(s) \bigr) \Bigr) = \mathbb EW^2(s) + \mathbb EW(s) \bigl( W(t) - W(s) \bigr) = \\
        = \mathbb EW^2(s) + \mathbb EW(s) \mathbb E \bigl( W(t) - W(s) \bigr) = \mathbb EW^2(s) = s
    \end{multline*}
    Аналогично можно рассмотреть ситуацию, когда $ t \le s $.

    В результате получаем, что $ \mathcal B(s, t) = \min\Set{s, t} $.
\end{undefthm}

\begin{property}
    Траектории винеровского процесса с вероятностью 1 \textbf{непрерывны}, но нигде \textbf{не
    дифференцируемы}.
\end{property}

Пусть $ H(x) = P\Set{\max\limits_{0 \le t \le 1} W(t) < x} $ "--- функция распределения
максимального значения того уровня, который достигается траекторией на временном интервале
$ [0, 1] $.
Для этой функции справедливо:
$$ H(x) =
\begin{cases}
    0, \quad x < 0 \\
    2\Phi(x) - 1, \quad x \ge 0
\end{cases} $$
